\documentclass[
    12pt,
    openright,
    oneside,
    a4paper,
    brazil
]{abntex2}

% Pacotes
\usepackage[utf8]{inputenc}
\usepackage[T1]{fontenc}
\usepackage{lmodern}
\usepackage{graphicx}
\usepackage{url}
\usepackage{microtype}
\usepackage{float}
\usepackage{csvsimple}
\usepackage{hyperref}

% Configuração de links azuis
\hypersetup{
    colorlinks=true,
    linkcolor=blue,
    citecolor=blue,
    urlcolor=blue
}

% Informações do trabalho
\titulo{Rede Meteorológica com Regador Inteligente}
\autor{Henry Thomaz Oliveira Campos}
\instituicao{
  Programa de Iniciação Científica e Tecnológica do Estado de Mato Grosso do Sul 2025
}
\local{Brasil}
\data{2025}
\orientador{Carlos Roberto Santos Junior}

\begin{document}

% Capa
\imprimircapa

% Folha de rosto
\imprimirfolhaderosto*

% Sumário
\renewcommand{\contentsname}{SUMÁRIO}
\tableofcontents
\cleardoublepage

% Conteúdo

\section{Introdução}
Este projeto tem como objetivo desenvolver uma rede de estações meteorológicas 
domésticas, capazes de monitorar dados climáticos (temperatura, umidade do ar, umidade 
do solo, chuva) e atuar de forma automática no acionamento de um regador inteligente. 
A ideia é evoluir em fases, começando com um protótipo funcional em uma única residência e 
chegando até uma rede colaborativa de estações interligadas com previsão via Inteligência 
Artificial.

\section{Justificativa}
O projeto busca resolver problemas práticos como:
\begin{itemize}
  \item Uso consciente de água na irrigação;
  \item Monitoramento climático acessível e comunitário;
  \item Educação tecnológica (eletrônica, IoT, programação e IA).
\end{itemize}

Além disso, se alinha com sustentabilidade e cidades inteligentes, dois temas de grande 
relevância.

\section{Objetivos}

\subsection{Geral}
Criar uma rede de estações meteorológicas inteligentes interligadas, com irrigação 
automatizada e integração de IA.  

\subsection{Específicos}
\begin{enumerate}
  \item Desenvolver protótipo inicial com ESP32, sensores e bomba d’água;
  \item Registrar dados meteorológicos em nuvem (Firebase/Supabase);
  \item Criar um painel web interativo com mapas e gráficos;
  \item Evoluir para rede colaborativa multiusuário;
  \item Implementar previsão climática com Inteligência Artificial.
\end{enumerate}

\section{Etapas do Projeto}

\subsection{Fase 1 – Protótipo básico (MVP)}  
Montagem com ESP32, sensores de solo e ar, bomba controlada via relé.  
Regras simples: “se solo seco → liga rega / se chuva → não rega”.  
Envio de dados para servidor gratuito.  
Site básico para visualização.  

\subsection{Fase 2 – Dados históricos e mapa}  
Armazenamento em banco de dados na nuvem.  
Visualização em mapa interativo com OpenStreetMap + Leaflet.  
Gráficos históricos com Chart.js.  

\subsection{Fase 3 – Rede de estações (multi-casa)}  
Cadastro de usuários e estações.  
Compartilhamento de dados entre vizinhos, familiares e comunidade.  

\subsection{Fase 4 – Inteligência Artificial}  
Treinamento de modelos de regressão/IA para prever variáveis climáticas.  
Exibição de dados reais + previstos no painel.  

\section{Protótipo Inicial – Estação Meteorológica}
Este é um protótipo inicial do projeto (Redes Meteorológicas). Ele é a base do 
projeto, mas muitas coisas mudarão, como o Arduino, que será substituído por um ESP32; o 
botão, que se tornará um sensor de chuva mais sofisticado; o potenciômetro, que será 
substituído por um sensor de velocidade do vento; o inversor, relés e o motor DC, que 
serão substituídos por uma bomba d’água.  

\vspace{0.5cm}

\textbf{Link do Protótipo:} 
\href{https://www.tinkercad.com/things/6btDA0XZOrf-estacao-meteorologica}{link-projeto}

\medskip

\textbf{Esquema do Protótipo:}
\begin{figure}[H]
    \centering
    \includegraphics[width=0.9\linewidth]{Estação Meteorológica_page-0001.jpg}
    \caption{Esquema do Protótipo}
    \label{fig:protótipo}
\end{figure}

\medskip

\textbf{Lista de Componentes (BOM):}

\begin{table}[H]
\centering
\begin{tabular}{|c|c|c|}
\hline
Código & Quantidade & Componente \\
\hline
\csvreader[head to column names]{bom.csv}{}%
{\csvcoli & \csvcolii & \csvcoliii \\ \hline}
\end{tabular}
\caption{Lista de Componentes do Protótipo (BOM)}
\label{tab:bom}
\end{table}

\section{Impacto Educacional e Social}
\begin{itemize}
  \item Promover aprendizagem em eletrônica, programação e ciência de dados;
  \item Incentivar o uso consciente da água;
  \item Pode ser expandido para escolas e comunidades como ferramenta de ensino;
  \item Usa cada estação para mesclar os dados em um mapa como diagramas de Venn e 
        fazer previsões meteorológicas mais avançadas.
\end{itemize}

\section{Conclusão}
O projeto vai além de um simples regador: é uma iniciativa de ciência cidadã, que conecta 
tecnologia, sustentabilidade e comunidade. Com um protótipo já iniciado e planejamento 
estruturado, o projeto tem grande potencial.

\begin{thebibliography}{99}

\bibitem{tinkercad} TINKERCAD. Protótipo base no Tinkercad. Disponível em: \href{https://www.tinkercad.com/things/3Hd7QsF61PI-regador-inteligente}{link-projeto}. Acesso em: 20 set. 2025.

\bibitem{firebase} GOOGLE. Firebase Documentation. Disponível em: \href{https://firebase.google.com/docs}{Firebase}. Acesso em: 20 set. 2025.

\bibitem{leaflet} LEAFLET. Leaflet.js. Disponível em: \href{https://leafletjs.com/}{Leaflet.js}. Acesso em: 20 set. 2025.

\bibitem{chartjs} CHART.JS. Chart.js Documentation. Disponível em: \href{https://www.chartjs.org/}{Chart.js}. Acesso em: 20 set. 2025.

\bibitem{manualmundo} MANUAL DO MUNDO. Disponível em: \href{https://youtu.be/_xRyePvaMqU?si=kJ5MQaMxgp4OxUS8}{vídeo}. Acesso em: 20 set. 2025.

\end{thebibliography}

\end{document}
